\chapter{Discussion and Conclusion}\label{chapter:closing}

\section{Limitations}\label{sec:limitations}

The headline numbers come from a single host running a multi-topic prompt at a fixed seed across both models.
Three caveats follow from that.

\begin{itemize}
\item The hardware is one consumer-class configuration: an 8-core AMD CPU, 30~GiB of DDR5, and a single PCIe~4.0 NVMe SSD as the model store.
A different CPU would shift the 22~GiB CPU-bound ceiling on GPT-OSS-20B.
A slower SSD would shrink the absolute speedups of the buffered modes.

\item Wall-clock numbers come from one prompt domain.
We separately captured tensor-level traces on five additional prompts (code, mathematical reasoning, creative writing, factual explanation, and a mixed code-plus-prose prompt).
The simulator's hit rate at the headline 250-slot LFU-aging configuration spans 32.6\% (mixed) to 43.0\% (creative), a 10.4 pp spread, and reuse-distance-1 spans 43--50\% across the same six domains.
Translating the hit-rate spread through our empirical hit-rate-to-tok/s sensitivity points to a few percent of cross-domain wall-clock variation, but we did not run end-to-end wall-clock measurements on the five extra domains.

\item Two replacement policies (ARC and W-TinyLFU) appear only in the simulator, not in the C runtime.
The simulator's predictions match the deployed policies to within rounding, but a live C implementation would close the loop in the operating regime where LRU and LFU-aging already perform within noise of each other.
\end{itemize}


\section{Future work}\label{sec:future-work}

Four directions extend naturally from where this thesis stops.

\paragraph{Cross-layer expert prefetching.}
The pipelining work in this thesis stops at the layer boundary, because the router's choice of next-layer experts is independent of the current layer's hidden state in any way we can predict cheaply.
A small router-prediction head, trained to forecast layer~$N{+}1$'s expert IDs from layer~$N$'s output, would let the loader start the disk reads one layer earlier and hide more of the I/O behind compute.
This is a separate research direction (expert speculation).

\paragraph{Request-similarity prefetch.}
The current system carries no state between requests.
Every request starts with an empty cache and pays the cold-start miss for each first-touch expert.
MoE-Infinity~\cite{xue2024moeinfinity} records each completed request as a small Expert Activation Matrix (rows are layers, columns are experts, entries are how often each expert was selected) and matches an in-flight request against historical ones by cosine similarity to drive prefetch.
The construction transfers cleanly to our setting: a few kilobytes per stored matrix, microsecond cosine matching, and a populated cache before the first decode token.
It targets a phase the per-token working-set arithmetic in \autoref{chapter:tracer} cannot address, and it is orthogonal to the router-prediction head above (cross-request fingerprint at request start versus cross-layer router prediction during decode).

\paragraph{Faster startup via lazy pinning.}
The current pinning mlocks the dense weight pages eagerly at model-load time.
Linux's \texttt{mlock2} with the \texttt{MLOCK\_ONFAULT} flag~\cite{mlock2_man} marks pages non-reclaimable but defers the actual residency until first fault.
For a generation-time workload this would shift the load-time cost into the first forward pass without changing steady-state behaviour.

\paragraph{GPU integration.}
The work in this thesis is entirely CPU-side.
The I/O architecture (the buffer manager, the sidecar pack file, the per-expert tagged completions) is orthogonal to the compute backend.
A GPU compute path would still need expert weights resident in some buffer before launch and the same routing-driven cache.
The integration question is the boundary: \texttt{O\_DIRECT} reads can target GPU-resident memory through GPUDirect Storage~\cite{nvidia_gds} or analogous facilities, pushing the buffer onto the device and keeping the CPU out of the data path.

\section{Conclusion}\label{sec:thesis-conclusion}

Modern open-source language models hold large amounts of weight data on disk.
The two studied here, GPT-OSS-20B and GPT-OSS-120B, weigh 12.85~GiB and 60.88~GiB respectively, against the 30~GiB of RAM on our test machine.
RAM is expensive, NVMe SSDs are not, and a Mixture-of-Experts model reads only a small fraction of its weights per generated token.

We built a tensor-level tracer to see what the model actually reads, and the trace showed a non-uniform access pattern: the dense per-layer weights are read on every token, while only a few experts per layer per token are selected by the router.
We pinned the dense weights in RAM so the kernel could not evict them, and we managed the experts ourselves in a buffer that reads from the SSD asynchronously, with three loading strategies of increasing aggressiveness.
The most aggressive submits all of a layer's expert reads at the start of the layer and dispatches each expert's compute as its data arrives.

With 7~GiB of RAM available, our backend lifts GPT-OSS-20B from 1.44~tok/s under the default loader to 5.36~tok/s, almost four times faster.
GPT-OSS-120B does not fit in our 30~GiB of RAM at all.
The default loader runs it at 4.96~tok/s and our backend at 8.06~tok/s.
A simulator built on the same trace lays a foundation for choosing the cache replacement policy and confirms the wall-clock observation that LFU-aging dominates when the cache is below the per-token working set, while LRU wins above it.

The tracer, the buffer-managed loader, the simulator, and the patched llama.cpp fork are released as artifacts (\autoref{chapter:artifacts}) so the path from question to answer can be walked again on different hardware, a different model, or a different storage stack.
